\begin{tikzpicture}[x=1cm,y=1cm,>=Stealth]
\tikzset{
  flowStart/.style={rectangle, rounded corners=1.5mm, draw, line width=0.6pt, minimum width=26mm, minimum height=10mm, inner xsep=3mm, inner ysep=1.5mm, align=center, execute at begin node={\setlength{\parindent}{0pt}\centering}},
  flowIO/.style={trapezium, trapezium left angle=75, trapezium right angle=105, draw, line width=0.6pt, minimum width=28mm, minimum height=10mm, inner sep=2mm, align=center, execute at begin node={\setlength{\parindent}{0pt}\centering}},
  flowProc/.style={rectangle, draw, line width=0.6pt, minimum width=32mm, minimum height=10mm, inner sep=2mm, align=center, execute at begin node={\setlength{\parindent}{0pt}\centering}},
  flowWideA/.style={rectangle, draw, line width=0.6pt, minimum width=58mm, minimum height=13mm, inner sep=2mm, align=center, execute at begin node={\setlength{\parindent}{0pt}\centering}},
  flowWideB/.style={rectangle, draw, line width=0.6pt, minimum width=58mm, minimum height=13mm, inner sep=2mm, align=center, execute at begin node={\setlength{\parindent}{0pt}\centering}},
  flowDec/.style={diamond, aspect=2.6, draw, line width=0.6pt, minimum width=36mm, minimum height=15mm, inner sep=1mm, align=center, execute at begin node={\setlength{\parindent}{0pt}\centering}},
  flowEdge/.style={->, >=Stealth, line width=0.6pt},
  flowPlain/.style={-, line width=0.6pt},
  flowLabel/.style={font=\normalsize, inner sep=1pt}
}
\node[flowStart] (start) at (0,0) {开始};
\node[flowIO] (input) at (0,-1.9) {输入 $t$};
\node[flowWideA] (init) at (0,-3.7) {$S=1,n=0,m=\dfrac{1}{2}$};
\coordinate (J) at (0,-5.0);
\node[flowProc] (sub) at (0,-6.2) {$S=S-m$};
\node[flowWideB] (upd) at (0,-7.9) {$m=\dfrac{m}{2},n=n+1$};
\node[flowDec] (dec) at (0,-9.9) {$S>t$};
\node[flowIO] (out) at (0,-11.9) {输出 $n$};
\node[flowStart] (end) at (0,-13.6) {结束};
\draw[flowEdge] (start.south) -- (input.north);
\draw[flowEdge] (input.south) -- (init.north);
\draw[flowEdge] (init.south) -- (sub.north);
\draw[flowEdge] (sub.south) -- (upd.north);
\draw[flowEdge] (upd.south) -- (dec.north);
\draw[flowEdge] (dec.south) -- (out.north) node[midway, right=2mm, flowLabel] {否};
\draw[flowEdge] (out.south) -- (end.north);
% 是回环：从判定框西顶点向左绕过两处理框，经外侧通道汇入主干J；标签置于外侧垂直段左侧以避开宽框
\draw[flowEdge] (dec.west) -- node[flowLabel, above=2mm, pos=0.22] {是} (-4.2,-9.9) -- (-4.2,-5.0) -- (J);
% label moved to first horizontal segment above (see draw above)
\end{tikzpicture}
